\subsection{Lazy Stream Evaluation: Generators}

A generator function contains \texttt{yield}. Calling it returns a generator object without running the body. Each \texttt{next()} or \texttt{send()} resumes the suspended \texttt{PyFrameObject} until the next \texttt{yield}, \texttt{return}, or exhaustion.

\subsubsection{The \texttt{yield} Keyword Architecture: Pausing Function Execution Without Destroying Local State}

%<<<PROGRAM:programs/laz02>>>


\texttt{yield} sends a value outward and suspends. Local variables (\texttt{x}, \texttt{y}, loop counters) remain in the generator frame.

\subsubsection{Generator Objects: Suspended Frames, Instruction Pointers, and Resumable Execution State}

A generator holds \texttt{gi\_code} (the \texttt{PyCodeObject}) and \texttt{gi\_frame} (the heap-allocated \texttt{PyFrameObject} while suspended). \texttt{f\_lasti} records the bytecode offset for resumption. When the generator completes, \texttt{gi\_frame} becomes \texttt{None} and the frame is deallocated.

%<<<PROGRAM:programs/laz01>>>


This confirms the general frame model: CPython keeps evaluation state on the heap so generators can pause without copying locals to a separate structure.

\subsubsection{Execution State Resumption: Re-Entering Suspended Function Frames Across Iteration Steps}

Each \texttt{next(gen)} re-enters the same frame at \texttt{f\_lasti}. The for-loop protocol calls \texttt{\_\_next\_\_} repeatedly until \texttt{StopIteration}.

\subsubsection{Returning from Generators: \texttt{StopIteration} and Generator Completion Values}

\texttt{return value} inside a generator raises \texttt{StopIteration} with \texttt{value} as the completion value (Python 3.3+). Callers using \texttt{next()} see the exception; \texttt{yield from} captures the completion value.
